\subsection{Feasibility}
\vspace{12pt}

The proposed project to develop and benchmark a modular ZK-Rollup framework for improving Ethereum's scalability is both technically and practically feasible, given the resources, timeline, and existing technologies. The feasibility is assessed across technical, resource, and time dimensions, with considerations for potential risks and mitigation strategies.

\subsubsection{Technical Feasibility}
\vspace{12pt}
The project leverages established tools and frameworks, such as zkSync and Polygon zkEVM, for core ZK-Rollup components, ensuring a robust foundation. The implementation uses well-documented technologies:
\vspace{12pt}
\begin{itemize}
    \item \textbf{Smart Contracts}: Solidity is a mature language for Ethereum, with libraries like OpenZeppelin providing reusable components for \texttt{RollupBridge.sol} and \texttt{Verifier.sol}.
    \item \textbf{Off-chain Sequencer}: Node.js/TypeScript is widely used for scalable off-chain systems, and Sparse Merkle Trees are standard for state management in ZK-Rollups.
    \item \textbf{Proof Systems}: Circom with Groth16/Plonk is production-ready for zk-SNARKs, with Plonky2 and Halo2 as viable alternatives for comparative experiments. Existing libraries reduce development overhead.
    \item \textbf{Data Availability}: EIP-4844 (proto-danksharding) is integrated into Ethereum testnets (e.g., Sepolia), and off-chain solutions like IPFS are accessible for prototyping.
\end{itemize}
\vspace{12pt}
The modular architecture ensures flexibility, allowing integration of different proof systems or transaction types (e.g., ERC20, ERC721). Deployment on an Ethereum testnet is straightforward, as testnets like Sepolia or Goerli provide free resources for experimentation. The use of synthetic workload generators and standardized benchmarking scripts aligns with industry practices, ensuring reliable metrics collection.
\vspace{12pt}

\subsubsection{Resource Availability}
\vspace{12pt}
The project is designed for a final-year research, requiring resources accessible to a student:
\vspace{12pt}
\begin{itemize}
    \item \textbf{Computational Resources}: 
    \vspace{12pt}
    \begin{table}[h!]
\centering
\begin{tabular}{lllll}
\toprule
\textbf{Phase} & \textbf{CPU} & \textbf{RAM} & \textbf{GPU} & \textbf{Storage} \\
\midrule
Dev/Testing & 4 cores & 16 GB & None & 250 GB SSD \\
Prototype Bench & 8-16 cores & 32-64 GB & RTX 3080 / A6000 & 1 TB SSD \\
\bottomrule
\end{tabular}
\caption{System requirements for different phases of ZK-Rollup prototype development and benchmarking.}
\end{table}
\vspace{12pt}


    \item \textbf{Software Tools}: Open-source tools (Circom, Hardhat, Node.js, Docker) are freely available. Benchmarking scripts and Jupyter notebooks for analysis require no additional licensing costs.
    \item \textbf{Knowledge and Skills}: The project assumes proficiency in Solidity, JavaScript/TypeScript, and basic cryptography (e.g., zk-SNARKs). Online resources, such as zkSync documentation and Ethereum developer guides, provide sufficient support for learning and troubleshooting.
\end{itemize}
\vspace{12pt}
Access to academic resources (e.g., IEEE papers, arXiv) and open-source communities (e.g., GitHub, Ethereum forums) further supports implementation and validation. No specialized hardware (e.g., GPUs for proof generation) is required for the prototype scale, as Circom and Groth16 are optimized for standard hardware.
\vspace{12pt}

\subsubsection{Time Feasibility}
\vspace{12pt}
The proposed timeline (July 2025 to April 2026, approximately 10 months) is realistic for the project. Key tasks, architecture design, prototype implementation, benchmarking, optimization experiments, analysis, and reporting are structured sequentially with buffers for delays. The modular design allows parallel work (e.g., data availability integration alongside testnet deployment) to optimize time. Potential delays, such as debugging proof generation or testnet instability, are mitigated by using established frameworks and allocating extra time in critical phases like implementation and experimentation.
\vspace{12pt}

\subsubsection{Risks and Mitigation}
\vspace{12pt}
Potential risks include:
\vspace{12pt}
\begin{itemize}
    \item \textbf{Complexity of ZK Proofs}: Circuit design and proof generation may be time-intensive. Mitigation: Start with simple circuits (e.g., token transfers) and use pre-built libraries like Circom.
    \item \textbf{Testnet Issues}: Network congestion or testnet upgrades could disrupt deployment. Mitigation: Use multiple testnets (e.g., Sepolia, Goerli) and local test environments (e.g., Hardhat).
    \item \textbf{Data Analysis Challenges}: Statistical analysis may reveal unexpected results. Mitigation: Use robust tools (e.g., Python with pandas, Jupyter) and repeat experiments for statistical significance.
\end{itemize}
\vspace{12pt}
The project’s focus on modularity and reproducibility ensures that partial outcomes (e.g., a working prototype or baseline benchmarks) are valuable even if experiments face delays.
\vspace{12pt}

\subsubsection{Societal and Economic Feasibility}
\vspace{12pt}
There is a significant and growing demand for scalable blockchain solutions across industries. The forecasted business value of blockchain (over \$3.1 trillion by 2030) and the high percentage of companies planning to integrate it underscore the economic incentive. Addressing scalability is crucial for blockchain to transition from experimental stages to mainstream production, enabling adoption in decentralised finance (DeFi), supply chain, healthcare, and IoT. This validates both the societal and economic relevance of the research.
\vspace{12pt}

\subsection{Limitations}
\vspace{12pt}
This project has several limitations that should be acknowledged. The prototype will be a proof-of-concept and not a production-ready system. It will only support a limited set of simple transactions (token transfers, deposits, withdrawals) and will not be EVM-compatible, which means it cannot run general-purpose smart contracts. The use of synthetic workloads, while based on real-world traces, may not capture the full complexity and unpredictability of live network traffic. Finally, while we will use a local testnet for controlled experiments, the benchmarks on the public testnet will be subject to network congestion and other external factors that are beyond our control.
\vspace{12pt}

\subsection{Conclusion}
\vspace{12pt}
This project proposes the design, implementation, and evaluation of a modular ZK-Rollup prototype on an Ethereum testnet, targeting scalable token transfers with configurable batching, sequencing, and data availability strategies. By combining an off-chain sequencer, state manager, prover module, and on-chain verification, the prototype provides a research-ready framework capable of exploring trade-offs between throughput, latency, gas costs, and storage. Through systematic benchmarking and optimization experiments, including variations in batch size, frequency, sequencing policies, proof systems, and data availability compression, the project aims to generate empirical evidence validating or challenging theoretical claims about ZK-Rollup scalability. The outcomes will provide a clear understanding of the practical performance and limitations of current ZK-Rollup designs. Moreover, the modular architecture ensures extensibility, allowing future experimentation with more complex transaction types, multi-sequencer setups, or production-grade zkEVM implementations. By open-sourcing our framework and findings, we hope to provide a valuable resource for the ZK-Rollup community and to inspire further research into the practical challenges of blockchain scalability.
\vspace{12pt}